\chapter*{常见缩写和术语先导}
\addcontentsline{toc}{chapter}{常见缩写和术语先导}

刷题材料里会频繁出现英文缩写，也会出现一些英文工程词。第一次见到它们时，不要先背模板，而要先知道它们在题目里负责表达什么对象、什么状态或什么证明边界。

\begin{description}
  \item[AC] Accepted，表示提交通过全部测试。它只能说明当前平台测试通过，不等于复杂度证明、边界证明和代码可维护性都已经完成。
  \item[LeetCode/LCP/LCR] LeetCode 是主站题库；LCP 通常指力扣杯竞赛题前缀；LCR 通常指力扣中国站重新编号的经典面试题集合。前缀不同只说明题库来源不同，读题、建模和证明方法不变。
  \item[II/III/IV] 题名后的罗马数字，表示同一主题的第二、第三、第四道变体，例如“课程表 II”。它不是新算法名，而是提醒你题目要求可能从判断变成构造、从一维变成多维。
  \item[reference/oracle] reference 是参考实现，通常慢但语义清楚，用来定义正确结果；oracle 是判定器，用同一批输入比较 reference 和优化版是否一致。刷题时先写 reference，可以避免直接背模板后不知道错在哪里。
  \item[GTest/GMock] GTest 是 GoogleTest，C++ 单元测试框架；GMock 是 GoogleMock，用来构造可控依赖和交互断言。本书实践工程目前用轻量测试主程序，但术语会在 C++ 工程和代码卷里出现，含义都是“用测试固定行为边界”。
  \item[C++20/C++23] C++ 标准版本名。C++20 表示本书代码默认使用的语言能力边界；C++23 是更新标准，刷题平台不一定支持，不能默认拿来写提交代码。
  \item[API] Application Programming Interface，接口。刷题里常指标准库函数或题目要求你实现的类方法，例如 \code{LRUCache::get} 和 \code{put}。
  \item[\code{std::}] C++ 标准库命名空间前缀。\code{std::vector}、\code{std::string}、\code{std::queue}、\code{std::priority_queue} 不是算法名，而是本书代码用来表达数组、字符串、队列和堆的标准库工具。
  \item[\code{int64\_t/size\_t}] \code{std::int64_t} 是固定 64 位有符号整数，常用于防止求和或乘法溢出；\code{std::size_t} 是容器大小和下标相关的无符号类型。刷题时看到它们，先问“值域是否会超过 int”和“下标是否可能为负”。
  \item[ID/IP] ID 是 identifier，表示身份或编号；IP 在“复原 IP 地址”中指 Internet Protocol 地址，题目把它简化成四段 0 到 255 的数字字符串切分问题。
  \item[I/O] Input/Output，输入输出。刷题里常见于“读入、输出”或工程题里的文件、网络边界；算法分析时通常不把平台 I/O 当作核心复杂度。
  \item[CPU] Central Processing Unit，中央处理器。刷题里提到 CPU 时，通常是在提醒时间复杂度最终会落到机器执行成本上，而不是让你研究硬件细节。
  \item[ASCII/UTF-8/Unicode] ASCII 是早期单字节字符编码集合；Unicode 是统一字符集合；UTF-8 是 Unicode 的常见字节编码方式。字符串题若只说小写英文，数组计数即可；若涉及 UTF-8，就不能把“字节数”直接当“字符数”。
  \item[STL] Standard Template Library，C++ 标准库中容器、迭代器和算法的常见统称。本书说 STL 时，重点是接口语义和复杂度，不是让你死记函数名。
  \item[INF/EPSILON] INF 表示人为设置的“足够大”哨兵值，常用于最短路和 DP 初始化；EPSILON 表示浮点比较容忍误差，避免把舍入误差当成逻辑差异。
  \item[DP] Dynamic Programming，动态规划。它不是“套表格”，而是把会重复遇到的子问题命名成状态，并保存状态答案，避免在搜索树中反复计算。
  \item[BFS] Breadth-First Search，广度优先搜索。它按层扩展状态，常用于无权最短路、最少步数和层序遍历。
  \item[DFS] Depth-First Search，深度优先搜索。它沿一条路径走到底再回退，常用于连通块、路径枚举、回溯和拓扑检测。
  \item[BST] Binary Search Tree，二叉搜索树。它要求左子树值小于当前节点、右子树值大于当前节点，因此中序遍历得到有序序列。
  \item[AVL] 一类自平衡二叉搜索树。它通过旋转控制左右子树高度差，让查找、插入和删除保持对数级。
  \item[DAG] Directed Acyclic Graph，有向无环图。只要状态依赖没有环，就可以用拓扑顺序或记忆化搜索求解。
  \item[SCC] Strongly Connected Component，强连通分量。有向图中一组点两两互相可达，常用 Tarjan 或 Kosaraju 算法求出。
  \item[LCA] Lowest Common Ancestor，最近公共祖先。树上两个节点往根方向走，第一次汇合的节点就是 LCA。
  \item[LRU] Least Recently Used，最近最少使用淘汰。核心是时间顺序，通常用哈希表定位节点、双向链表维护新旧顺序。
  \item[LFU] Least Frequently Used，最不经常使用淘汰。它比 LRU 多一个频次维度：频次相同时再按时间淘汰。
  \item[LCS/LIS] LCS 是最长公共子序列，LIS 是最长递增子序列。二者都属于序列 DP，但状态含义不同，不能混用模板。
  \item[KMP] Knuth-Morris-Pratt 字符串匹配算法。它用失败函数记录“当前匹配失败后还能保留多长前后缀”，避免主串指针回退。
  \item[KM] Kuhn-Munkres 算法，常用于二分图最大权完美匹配。它和无权最大匹配的普通匈牙利算法不是同一个问题边界。
  \item[DAWG] Directed Acyclic Word Graph，有向无环词图，可看作压缩后的字符串自动机/字典结构，用更少状态表达大量后缀或单词关系。
  \item[SPFA] Shortest Path Faster Algorithm，Bellman-Ford 的队列优化写法。它能处理负权边，但复杂度不稳定，面试中要谨慎说明适用边界。
  \item[TSan] ThreadSanitizer，线程检查器。它能帮助发现数据竞争，但不能替代并发算法的不变量证明；刷题书里主要在 C++ 工程化和并发题背景中出现。
  \item[ABA] ABA problem，无锁结构里的经典问题：共享值从 A 变成 B 又变回 A，观察者只看当前值会误以为没有变化。链表、栈、RCU 和内核结构章节提到它时，重点是对象身份和生命周期不能只靠表面值判断。
  \item[IDA*] Iterative Deepening A*，带启发式估价的迭代加深搜索。它常用于状态空间很大但答案深度可逐步放宽的问题。
  \item[2-SAT] 2-satisfiability，二元布尔可满足性问题。刷题里通常把“选 A 就不能选 B”建成蕴含图，再用强连通分量判断矛盾。
  \item[mask/bitmask] mask 是掩码，bitmask 是用二进制位保存集合或状态。看到 \code{1 << i} 时，要先问每一位代表哪个选择、状态如何合并。
  \item[URBG] Uniform Random Bit Generator，C++ 随机库里的“均匀随机比特生成器”概念。看到它时先理解成“能持续产生随机数的引擎接口”。
  \item[DNA] Deoxyribonucleic Acid，脱氧核糖核酸。字符串题里常把 \code{A/C/G/T} 当作四种字符，重点仍是哈希、滑窗或位编码。
  \item[IPO] Initial Public Offering，首次公开募股。对应题目通常把项目利润和启动资本抽象成“先选可启动项目，再最大化资本”的贪心/堆模型。
  \item[URL] Uniform Resource Locator，统一资源定位符。字符串题中出现 URL 时，通常不是考网络协议，而是考字符转义、路径切分或状态机解析。
  \item[RCU] Read-Copy-Update，Linux 内核常见同步思想。读者尽量无锁读，更新者复制并替换对象，旧对象等读者离开后再回收。
  \item[SLAB/SLUB] Linux 内核的小对象分配器族，用缓存固定大小对象来降低频繁分配成本。SLUB 是现代 Linux 常见实现之一。
  \item[VMA] Virtual Memory Area，Linux 进程地址空间中的一段连续虚拟内存区域，例如一段 mmap 文件或堆区。
  \item[XArray] Linux 内核中的稀疏整数索引结构，常用于把整数 id 映射到对象。
  \item[CSR] Compressed Sparse Row，压缩稀疏行格式。刷题里它常用于图或稀疏矩阵，把每个点的邻居压进连续数组，遍历时从下标范围取邻边。
  \item[FIFO/LIFO] FIFO 是先进先出，典型结构是队列；LIFO 是后进先出，典型结构是栈。
  \item[target/head/tail] target 常表示题目要达到的目标值或目标状态；head/tail 常表示链表、队列、窗口或双指针的前后端。读到这些词时，要先把它们落回题目对象，而不是当作固定变量名。
  \item[\code{NEG\_INF/MAX\_VALUE}] 这类全大写名字通常是哨兵常量。NEG\_INF 表示“不可能小到的初值”，MAX\_VALUE 表示上界；使用前必须确认不会和真实答案混淆，也不会在加减时溢出。
  \item[\code{ALPHABET\_SIZE/ASCII\_RANGE}] 字符计数数组的容量常量。ALPHABET\_SIZE 可能只表示 26 个小写字母，ASCII\_RANGE 可能表示 128 或 256 个字节值；它们由题目字符集决定，不是固定模板。
  \item[样例字符串] \code{ABC}、\code{AABC}、\code{ADOBECODEBANC}、\code{BANC} 这类大写连续字母通常是题目样例里的字符串内容，不是缩写。读滑动窗口、子序列或回溯题时，先把它们当作具体输入逐步模拟。
  \item[命名常量] \code{DEFAULT_CAPACITY}、\code{GROWTH_FACTOR}、\code{SEGMENT_COUNT}、\code{ITERATIONS} 这类名字是代码为了消除魔法数字而取的常量名。它们的含义来自当前题目的容量、增长倍率、段数或循环次数，不是算法分类。
  \item[token/hash/partition/frontier] token 是被切出来的一个词元或片段；hash 是把对象映射成便于比较或查找的摘要；partition 是把数组、集合或状态空间切成几块；frontier 是 BFS 或图搜索当前这一层即将扩展的边界集合。
  \item[stack/kernel/continuation] stack 是栈，可以是数据结构，也可以是递归调用栈；kernel 在算法题里多指核心计算逻辑，不是操作系统内核；continuation 表示“后续还要执行的状态”，常出现在搜索、协程或异步题的解释里。
  \item[PDF] 固定版式电子书格式，适合打印、标注和精确页码。本书正式导出只保留 PDF。
  \item[代码状态常量] \code{DIRECTIONS}、\code{FRESH}、\code{ROTTEN}、\code{LAND}、\code{WATER}、\code{RED/WHITE/BLUE} 这类大写词通常不是新算法缩写，而是代码里给方向、状态或颜色取的名字。读到它们时，把它们还原成题目状态，而不是背字母。
  \item[C++ 容器操作] \code{push\_back} 是把元素追加到容器尾部；\code{emplace\_back} 是在容器尾部直接构造元素；\code{pop\_back} 和 \code{pop\_front} 分别从尾部、头部移除元素。它们是标准库接口动作，不是算法结论；看到它们时要问“这个结构维护的顺序、不变量和复杂度是什么”。
  \item[二分库函数] \code{lower\_bound} 返回第一个不小于目标值的位置，\code{upper\_bound} 返回第一个大于目标值的位置。它们要求搜索区间已经按比较规则有序；如果题目没有先保证有序，直接套函数就是错的。
  \item[常见模板工具] \code{static\_cast} 是 C++ 显式类型转换，通常用于把索引、计数或字符值转成目标类型；\code{numeric\_limits} 查询某个数值类型的最大值、最小值等边界。它们服务边界表达，不改变算法本身。
  \item[底层 C++ 工具] \code{shared\_ptr} 是引用计数智能指针，多个对象共享所有权；\code{reinterpret\_cast} 是低层位模式重解释，常见于讲内核式结构或地址布局的章节。算法题里一般不需要 reinterpret\_cast；出现它时，重点是理解对象身份、内存布局和生命周期风险。
  \item[题目结构名] \code{ListNode}、\code{TreeNode}、\code{LRUCache}、\code{DisjointSetUnion}、\code{FenwickTree} 这类 CamelCase 名字通常是题目或本书给出的结构类型。ListNode 表示链表节点，TreeNode 表示树节点，DisjointSetUnion 是并查集，FenwickTree 是树状数组；先理解结构维护什么关系，再看成员函数。
  \item[LeetCode 方法名] \code{sumRegion}、\code{maxFreq}、\code{findMedianSortedArrays} 这类名字来自题目要求实现的函数签名。它们不是新术语，真正要读的是函数输入、输出和必须满足的复杂度。
  \item[变量命名风格] \code{ending\_here}、\code{prefix\_sum}、\code{current\_end}、\code{next\_index} 这类 snake\_case 名字是为了让代码自解释：ending\_here 表示“以当前位置结尾的状态”，prefix\_sum 表示“到某个边界为止的累加和”，current\_end 表示“当前段的右边界”。读变量名时要把它放回当前题目的状态定义里，不要单独背名字。
  \item[复杂度短写] \code{V+E} 常表示图算法复杂度中的“顶点数加边数”；\code{K-1}、\code{K+1} 这类写法通常是围绕第 K 个位置或选择数做边界推导。读复杂度时，必须先确认每个字母代表题目里的哪个规模。
  \item[题目状态词] \code{ACTIVE}、\code{EMPTY}、\code{WILDCARD}、\code{TARGET}、\code{BASE} 这类全大写词多数是代码常量：ACTIVE 表示活跃，EMPTY 表示空位，WILDCARD 表示通配符，TARGET 表示目标，BASE 表示进制或基准值。它们的真实含义以当前题目建模为准。
\end{description}
